\documentclass[12pt, twoside]{article}
% Emit local LDF shims before packages load
\input{lib/datatool-ldf}

% Make it possible to conditionally depend on the TeX engine used
\RequirePackage{ifluatex}

\RequirePackage{comment}

% To access the \captionof{} functionality
\RequirePackage{caption}

%\RequirePackage{silence}
%\WarningFilter{latexfont}{Font shape `U/stmry/m/n' in size <5.5> not available}
%\WarningFilter{latexfont}{Font shape}

\newif\ifinswedish
\inswedishfalse

\RequirePackage{etoolbox}

% The following toggles are needed because of kth/kthcolors.tex and lib/includes.tex
\newif\ifnomenclature
\nomenclaturefalse

\newif\ifdigitaloutput
\digitaloutputfalse

\makeatletter
%% Define a pair of commands to disable and reenable specific packages - see https://tex.stackexchange.com/questions/39415/unload-a-latex-package
\makeatletter
\newcommand{\disablepackage}[2]{%
  \disable@package@load{#1}{#2}%
}
\newcommand{\reenablepackage}[1]{%
  \reenable@package@load{#1}%
}
\makeatother

%% To avoid the warning: "Package transparent Warning: Loading aborted, because pdfTeX is not running in PDF mode."
\disablepackage{transparent}{}

% include a variety of packages that are useful
\input{lib/includes}

\usepackage{minted}
\RequirePackage{fancyhdr} % Take control of headers and footers

%\usepackage{luatexja-fontspec}
%\ltjdefcharrange{8}{"2190-"21FF} % Define range 8 as Arrows
%\ltjsetparameter{jacharrange={-8}} % Treat range 8 as Western characters (ALChar)

%\ltjsetparameter{jacharrange={-3}} % Treat range 3 as Western characters (ALChar)

\begin{comment}
\ifinswedish
    \usepackage[english, main=swedish, bidi=basic]{babel}
\else
    \usepackage[swedish, main=english, provide+=*, bidi=basic]{babel}
\fi
\end{comment}
\ifinswedish
    \usepackage[main=swedish]{babel}
\else
    \usepackage[main=english, provide+=*]{babel}
\fi
%\setmainjfont{Noto Sans CJK JP}

%\usepackage[style=ieee,citestyle=numeric-comp, maxbibnames=99, giveninits=false]{biblatex}
% Input centralized biblatex configuration
\input{lib/biblatex-and-friends}



% Use the chapter* style for the heading of the bibliography
\defbibheading{bibliography}[\bibname]{%
\section*{#1}%   skip \centering 
}
  
\addbibresource{references.bib}    
\input{kth/kthcolors}



\input{lib/defines}  % load some additional definitions to make writing more consistent

\usepackage[plainpages=false, unicode=true]{hyperref}
\input{lib/includes-after-hyperref}
\usepackage{orcidlink}  %% to provide ORCID icons and links

%\usepackage[acronym, style=super, toc=false, nonumberlist, nomain, nopostdot=true, notranslate]{glossaries-extra}
% In the document preamble where glossaries-extra is loaded:
\usepackage[
  acronym,
  toc=false,
  nonumberlist,
  nomain,
  nopostdot=true, 
  nohypertypes={main, acronym}, % disables link generation to targets that do not exist
  notranslate,
]{glossaries-extra}
%% In overleaf, if this file is in a folder and not the root
%% using \makeglossaries will not work, as the Overleaf latexmk
%% will not run the program to take the output.acn file and make the output.acr file
%\makeglossaries
% However, you can get TeX to do the work with the following:
\makenoidxglossaries
% For details of the Overleaf make process, see https://www.overleaf.com/learn/how-to/How_does_Overleaf_compile_my_project%3F
%% Potentially, one could set up a latexmkrc file in the folder

\input{lib/acronyms}                %load the acronyms file

% Include Glossary ---
% Align the text expansion of the glossary entries
\newglossarystyle{mylong}{%
  \setglossarystyle{long}%
  \renewenvironment{theglossary}%
     {\begin{longtable}[l]{@{}p{\dimexpr 2cm-\tabcolsep}p{0.8\hsize}}}% <-- change the value here
     {\end{longtable}}%
 }
 
% define a left-aligned table cell that is ragged right
\newcolumntype{L}[1]{>{\raggedright\let\newline\\\arraybackslash\hspace{0pt}}p{#1}}

\usepackage{cleveref}           %% Replace Section with a symbol
\usepackage{metalogo}   % for \LuaLaTeX logo

  \RequirePackage{fontspec}
  \defaultfontfeatures{Ligatures={TeX}} % This enables TeX style ligatures such as ---, '', ``, and so on

    \babelfont{rm}{TeX Gyre Termes}
    \DeclareFontShape{TU}{TeXGyreTermes(0)}{md}{n}{<->sub * TeXGyreTermes(0)/m/n}{}
    \DeclareFontShape{TU}{TeXGyreTermes(0)}{sb}{n}{<->ssub * TeXGyreTermes(0)/b/n}{}
    \DeclareFontShape{TU}{TeXGyreTermes(0)}{sb}{it}{<->ssub * TeXGyreTermes(0)/b/it}{}
  
    %\setsansfont{TeX Gyre Heros}   %% Helvetica like font
    \babelfont{sf}{TeX Gyre Heros}
    %\setmonofont[Ligatures={NoCommon}, Numbers={Lining,Monospaced}]{TeX Gyre Cursor}  %% Courier like font
    %% Turn off ligature for tt font
    \babelfont{tt}[Ligatures={ResetAll}]{TeX Gyre Cursor}
 \begin{comment}
    \babelfont[english]{rm}{TeX Gyre Termes}
    \babelfont[english]{sf}{TeX Gyre Heros}
    \babelfont[english]{tt}{TeX Gyre Cursor}
\end{comment}

 %  \setmathfont{TeX Gyre Termes Math} %% a math font
    \usepackage{mathtools}
 \usepackage[warnings-off={mathtools-colon,mathtools-overbracket}]{unicode-math}
  % The [version=bold, FakeBold=1.2] is to avoid a warning about the lack of a bold font
  %\setmathfont{TeX Gyre Pagella Math}[version=bold, FakeBold=1.2] %% a font for math
  \setmathfont{STIX Two Math}[version=normal]
  %\setmathfont{TeX Gyre Pagella Math}[version=normal]
  %\setmathfont{TeX Gyre Pagella Math}[version=bold, BoldFeatures={FakeBold=1.5}]
  % For both XeLaTeX and LuaLatex for getting access to unicode symbols
  %\newfontfamily\myfont[CharacterVariant=1]{NewCM10-Regular.otf}
  % STIX Project (Scientific and Technical Information Exchange)
  % STIX Two Math does not have bold face - so we fake it
  %\newfontfamily\mystixmathfont[BoldFeatures={FakeBold=1.5}, BoldItalicFeatures={FakeBold=1.5}]{STIX Two Math}
  \newfontfamily\mystixmathfont{STIX Two Math}
  % STIX Two Math does not have a bold font, but it has bold symbols with and without serifs - but you manually have to use them, unless you are in math mode - then you can use \symbf{}
  % and this will return the bold serif version of the character
\makeatletter
\AtBeginDocument{%
  \@for\idx:={0,1,2,3,4,5,6}\do{%
    \@ifundefined{T@TU+STIXTwoMath(\idx)}{%
      % Family not loaded by unicode-math; do nothing
    }{%
      \DeclareFontShape{TU}{STIXTwoMath(\idx)}{b}{n}{<->ssub * STIXTwoMath(\idx)/m/n}{}%
      \DeclareFontShape{TU}{STIXTwoMath(\idx)}{sb}{n}{<->ssub * STIXTwoMath(\idx)/m/n}{}%
      \DeclareFontShape{TU}{STIXTwoMath(\idx)}{b}{it}{<->ssub * STIXTwoMath(\idx)/m/it}{}%
      \DeclareFontShape{TU}{STIXTwoMath(\idx)}{sb}{it}{<->ssub * STIXTwoMath(\idx)/m/it}{}%
    }%
  }%
}
\makeatother
  \newfontfamily\mystixtextfont{STIX Two Text}
  
    % To access the Stylistic Set 1 <ss01> such as ℋ
    \newfontfamily\mystixmathfontSSa{STIX Two Math}[StylisticSet=1]

    % use english as a fallback when in other languages
    \babelprovide[import, onchar=ids fonts]{english}

    
  % for new KTH cover
  % Load the Figtree font as it is used for the new KTH graphical profile
  % 

    \newfontfamily{\FigtreeFont}[Ligatures=TeX,
        Path=./Figtree/static/,
        Extension = .ttf,
        UprightFont=*-Regular,
        BoldFont=*-Bold,
        BoldItalicFont=*-BoldItalic,
        ItalicFont=*-Italic,
        %FontFace={l}{n}{*-Light},
        %FontFace={l}{it}{*-LightItalic},
        FontFace={md}{n}{*-Medium},
        FontFace={md}{it}{*-MediumItalic},
        FontFace={sb}{n}{*-Semibold},
        FontFace={sb}{it}{*-SemiBoldItalic},
        %FontFace={k}{n}{*-Black},
        %FontFace={k}{it}{*-BlackItalic},
        %FontFace={eb}{n}{Font=*-ExtraBold},
        %FontFace={eb}{it}{Font=*-ExtraBoldItalic}
        ]{Figtree}


\newfontfamily\pageNumberFont{Figtree} %% set the font to use for page numbering

\newfontfamily{\NotoEmojiFont}[Ligatures=TeX,
    Path=./Noto_Emoji/static/,
    Extension = .ttf,
    UprightFont=*-Regular,
    BoldFont=*-Bold,
    FontFace={l}{n}{*-Light.ttf},
    FontFace={md}{n}{*-Medium},
    FontFace={sb}{n}{*-SemiBold},
    ]{NotoEmoji}



\begin{comment}
   % To set the abstract headings in Figtree, we redefine the abstract environment to look at the language being used and use the appropriate font, with the default being Figtree
   % The languages that are automatically introduced by Polyglossia or Babel have a name of the form xxxfont and xxxfontsf; where xxxfont is the serif font and xxxfontsf is the sans serif font.
   % This means that for each language that Figtree does not support, you have to define the sans serif and serif font to use.

      \babelprovide[import, onchar=ids fonts]{greek}
      \babelprovide[import, onchar=ids fonts]{russian}
      \babelprovide[import, onchar=ids fonts]{vietnamese}
      \babelfont[greek, russian, vietnamese]{rm}{Noto Serif}
      \babelfont[greek, russian, vietnamese]{sf}{Noto Sans}
      \babelfont[greek, russian, vietnamese]{tt}{Noto Mono}

  
    \babelprovide[import, onchar=ids fonts]{hindi}
    \babelfont[hindi]{rm}{Noto Serif Devanagari}
    \babelfont[hindi]{sf}{Noto Sans Devanagari}
    \babelfont[hindi]{tt}{Noto Sans Devanagari} % Noto Mono does not have the glyphs

    %\babelprovide[import, onchar=ids fonts]{russian}
    %\babelfont[russian]{rm}{Noto Serif}
    %\babelfont[russian]{sf}{Noto Sans}
    %\babelfont[russian]{tt}{Noto Mono}

    \babelprovide[import, onchar=ids fonts]{chinese-simplified}
    \babelfont[chinese-simplified]{rm}{Noto Serif CJK SC}
    \babelfont[chinese-simplified]{sf}{Noto Sans CJK SC}
    \babelfont[chinese-simplified]{tt}{Noto Sans Mono CJK SC}
    
    \babelfont[chinese-traditional]{rm}{Noto Serif CJK TC}
    \babelfont[chinese-traditional]{sf}{Noto Sans CJK TC}
    \babelfont[chinese-traditional]{tt}{Noto Sans Mono CJK TC}
  
    \babelprovide[import, onchar=ids fonts]{japanese}
    \babelfont[japanese]{rm}{Noto Serif CJK JP}
    \babelfont[japanese]{sf}{Noto Sans CJK JP}
    \babelfont[japanese]{tt}{Noto Sans Mono CJK JP}
  
  % If you are going to use Arabic

    \babelprovide[import, onchar=ids fonts]{arabic}
    \babelprovide[import, onchar=ids fonts]{centralkurdish}
    \babelfont[arabic, centralkurdish]{rm}{Noto Naskh Arabic}
    \babelfont[arabic, centralkurdish]{sf}{Noto Sans Arabic}
    \babelfont[arabic, centralkurdish]{tt}{Noto Sans Arabic}
    % If one really needs a monospaced font, one might try Kawkab Mono
    % However, it seems that it is a work in progress - see https://makkuk.com/kawkab-mono/ and https://github.com/aiaf/kawkab-mono/tree/master

    %\babelprovide[import, onchar=ids fonts]{centralkurdish}
    %\babelfont[centralkurdish]{rm}{Noto Naskh Arabic}
    %\babelfont[centralkurdish]{sf}{Noto Sans Arabic}
    %\babelfont[centralkurdish]{tt}{Noto Sans Arabic}

      
  % If you are going to use Hebrew or Yiddish
\babelprovide[import, onchar=ids fonts]{hebrew}
\babelprovide[import, onchar=ids fonts]{yiddish}
\babelfont[hebrew, yiddish]{rm}{Noto Serif Hebrew}
\babelfont[hebrew, yiddish]{sf}{Noto Sans Hebrew}
\babelfont[hebrew, yiddish]{tt}{Noto Sans Hebrew}


    %\babelprovide[import, onchar=ids fonts]{vietnamese}
    %\babelfont[vietnamese]{rm}{Noto Serif}
    %\babelfont[vietnamese]{sf}{Noto Sans}
    %\babelfont[vietnamese]{tt}{Noto Mono}

\end{comment}
  
    % The Overleaf TeX Live includes these fonts, so there is little you have to do!
    % The list of such fonts is at https://www.overleaf.com/learn/latex/Questions%2FWhich_OTF_or_TTF_fonts_are_supported_via_fontspec%3F
    %
\newfontfamily{\NotoSansJPMonoFont}[Ligatures=TeX,
    ]{Noto Sans Mono CJK JP Regular}


\newfontfamily{\NotoSansJPFont}[Ligatures=TeX,
    ]{Noto Sans CJK JP}


\newfontfamily{\NotoSansFont}[Ligatures=TeX,
Extension = .ttf,
    UprightFont = NotoSans-Regular,
    BoldFont = NotoSans-Bold,
    ItalicFont = NotoSans-Italic,
    BoldItalicFont = NotoSans-BoldItalic
    ]{NotoSansCustomA}
    
\newfontfamily{\NotoSerifFont}[Ligatures=TeX,
Extension = .ttf,
    UprightFont = NotoSerif-Regular,
    BoldFont = NotoSerif-Bold,
    ItalicFont = NotoSerif-Italic,
    BoldItalicFont = NotoSerif-BoldItalic
    ]{NotoSerifCustomA}

\newfontfamily{\DejaVuSansFont}[Ligatures=TeX,]{DejaVu Sans}

% A font to be used in minted listings
\newfontfamily{\ListingMono}{Noto Sans Mono}[Ligatures=ResetAll]

% fonts for various symbols
\newfontfamily{\NotoSansSymbolsFont}{Noto Sans Symbols}[Ligatures=ResetAll]
\newfontfamily{\NotoSansSymbolsTwoFont}{Noto Sans Symbols 2}[Ligatures=ResetAll]



% color symbols - this would require changes in the font used in the U+2600-U+26FF-Miscellaneous-Symbols.tex file
% for example for characters: ☕, ⚽, ⛄, and ✈
%\newfontfamily{\NotoColorEmojiFont}{Noto Color Emoji}[Renderer=Harfbuzz]

% Phonetic Extensions, U+1D00 - U+1D7F
\input{unicode_blocks/U+1D00-U+1D7F-Phonetic-Extensions}
% Superscripts and Subscripts, U+2070 - U+209F
\input{unicode_blocks/U+2070-U+209F-Superscripts-and-Subscripts}
% Miscellaneous Symbols, U+2600 - U+26FF
\input{unicode_blocks/U+2600-U+26FF-Miscellaneous-Symbols}

% Set up the header and footer for plain style pages
\fancypagestyle{plain}{%
\fancyhf{}% clear all header and footer fields
\fancyfoot[C]{\pageNumberFont\small\selectfont\thepage}
\renewcommand{\headrulewidth}{0pt}%
\renewcommand{\footrulewidth}{0pt}%
}
% Set up the header and footer
\fancyhead{}
%\fancyhead[RO]{\sffamily\small\leftmark\thinspace|\thinspace\thepage}
%\fancyhead[LE]{\sffamily\small\thepage\thinspace|\thinspace\leftmark}
% Without conversion to uppercase
% Note that we need to switch to the familydefault for the page numbers
\fancyhead[RO]{{\sffamily\small\nouppercase\leftmark}{\pageNumberFont\small\selectfont \thinspace|\thinspace \thepage}}
\fancyhead[LE]{{\pageNumberFont\small\selectfont\thepage \thinspace|\thinspace}{\sffamily\small\nouppercase\leftmark}}
\fancyfoot{}
\renewcommand{\headrulewidth}{0pt}
\pagestyle{fancy}
\renewcommand{\sectionmark}[1]{%
\markboth{#1}{}}


\setlength{\headheight}{15pt}
\addtolength{\topmargin}{-3pt}

\input{lib/configure_listings}
% If you use both listing and lstlisting environments - the following makes them both use the same counter and the same formatting in the List of Listings
\makeatletter
\AtBeginDocument{\let\c@listing\c@lstlisting}
\AtBeginDocument{\let\l@listing\l@lstlisting}
\makeatother

\newcommand{\dname}[1]{\textbf{#1}}
\newcommand{\fname}[1]{\texttt{#1}}


\RequirePackage{luacode}
\input{kth/Lua_timestamping}

% Input file generators before \documentclass
\input{lib/ext-eprint-dbx}

\input{lib/check_for_placeholder_references}

\title{README and notes about the template for programmers --- Appendix}

\author{Gerald Q. Maguire Jr.}
\date{October 2026}

\begin{document}
\pagenumbering{roman}
\glsresetall
\maketitle
\fancypagestyle{empty}{}
\warningExpl{\textbf{This document is a work in progress.}}

\begin{abstract}
This document contains an appendix for use with \fname{README\_programmer\_notes.tex}.


\warningExpl{{\mystixmathfont ☡} This is a very deep dive that is probably only of interest to a \textit{very} small number of readers.}
\end{abstract}
\cleardoublepage

%%%%%%%%%%%%%%%%%%%%%%%%%%%%%%%%%%%%%%%%%%%%%%%%%%%%%%%%%%%%%%%%%%%%%%
% Enable the following command the first time you compile this in a new session to give all of the time quanta to set up the font cache.
%% If the font cache is not built, you are likely to get a timeout.
%\end{document}

% Include Table of Contents (TOC) ---
\fancypagestyle{plain}{}
\renewcommand{\sectionmark}[1]{ \markboth{#1}{}}
\tableofcontents
\markboth{\contentsname}{}
\cleardoublepage
 
\setcounter{page}{1}
\pagenumbering{arabic}
\glsresetall

This document provides a set of appendices for~\cite{maguire2026programmersguide_testing}.

\appendix
\section{Language codes}
\label{sec:LanguageCodes}
The US Library of Congress has a table of the language codes at \url{https://www.loc.gov/standards/iso639-2/php/code_list.php}.

Based upon the US Library of Congress \enquote{Codes for the Representation of Names of Languages: Codes arranged alphabetically by alpha-3/ISO 639-2 Code}, I have eliminated all of the rows that did not have both a three-letter \gls{ISO} 639-2 Code \textbf{and} a two-letter \gls{ISO} 639-1 Code. For the three-letter \gls{ISO} 639-2 Codes, I used only the  \enquote{B}, \ie bibliographic codes. I split the table into two parts: Languages already implemented (\Cref{tab:languagesImplemented}) and those remaining to be implemented (\Cref{tab:languagesNotImplemented}). The term \enquote{implemented} means that I have considered it in the current template, at least to some extent (but do not necessarily have a skeleton abstract in the language\footnote{For example, I do not have a skeleton abstract in Belarusian - but babel supports belarusian and the template supports the use of the Cyrillic alphabet. However, I have not done an investigation of whether the font supports all of the glyphs used in Belarusian.}).  For example, babel supports northernsami and yiddish as languages you can select, and I have created skeleton abstracts for both languages. In the case of Yiddish, I added it to the language switch in the code described in \textit{Extending the set of languages supported for abstracts} in \fname{README\_programmer\_notes}. For Norwegian Bokmål, one can use the \enquote{nob} code for an abstract and specify the babel language Norwegian. In contrast, for Norwegian Nynorsk one uses the \enquote{nno} code for an abstract and the babel language nynorsk. However, I have only implemented a skeleton abstract for Norwegian.

I have a skeleton abstract in Kurdish Sorani (babel language centralkurdish) with language code \enquote{cbk}, I do not have any support for Kurmanji. However, as Kurmanji uses a Latin-based alphabet, it should not be a problem to add an abstract with the \enquote{kur} code.

According to SCB (\url{https://www.scb.se/contentassets/221c69e399f0424b852b21236952e300/sprakkoder.xlsx}), there are also \gls{ISO} 639-2 codes for Yiddish with regions (see \Cref{tab:SCByiddish}). For Romanian, the \gls{ISO} 639-2 B code is \enquote{rum} and the ISO 639-1 code is  \enquote{ro}. However, SCB gives additional \gls{ISO} 639-2 codes (see \Cref{tab:SCBromani}). For Northern Sami (Samiska, (norra)), the \gls{ISO} 639-2 is "sme" and the ISO 639-1 code is \enquote{se}, but is there a region code that should be used? (See \Cref{tab:SCBsami}) For Tornedalen Finnish (\foreignlanguage{swedish}{Meänkieli}), the ISO 639-2 code is \enquote{fit}, but I am unable to determine what \gls{ISO} 639-1 code should be used (see \Cref{tab:languagesMissing}). Is it \enquote{fi} with some region code? If so, what region code? 

\generalExpl{I have added a skeleton abstract for Tornedalen Finnish (\foreignlanguage{swedish}{Meänkieli}). It is probably incorrect, but it provides a starting point. My assumption is that the spacing, hyphenation, \etc follow that of Finnish, but that there is a difference in words.}

\begin{table}[!htp]
    \caption{SCB table - Yiddish with regions}
    \label{tab:SCByiddish}
    \begin{tabular}{L{3cm}|L{6.5cm}}
      \textbf{\gls{ISO} 639-2 Code} &\textbf{Swedish name of Language}\\
      \hline
YIH  &  Jiddisch (Väst)\\
YDD  &  Jiddisch (Öst)\\
\end{tabular}
\end{table}

\begin{table}[!htp]
    \caption{SCB table - Romainian with regions}
    \label{tab:SCBromani}
    \begin{tabular}{L{3cm}|L{6.5cm}}
      \textbf{\gls{ISO} 639-2 Code} &\textbf{Swedish name of Language}\\
      \hline
ROM &   Romani chib\\
RMO &   Romani, Abbruzzesi, Serbisk romani, Slovensk-kroatisk romani\\
RMN &   Romani, Arli, Dzambasi, Gurbeti\\
RMC &   Romani, Bashaldo, Ungersk-slovakisk romani, Romungro\\
RMF &   Romani, Kale, Kalo\\
RMY &   Romani, Lovara, Kalderash\\
RML &   Romani, Polsk romani, Estnisk, Lettisk, Nordrysk, Vitrysk\\
RMU &   Romani, resande romani, Svensk romani, Romani tavringer\\
RMW &   Romani, Walesisk\\
\end{tabular}
\end{table}

\begin{table}[!htp]
    \caption{SCB table - Sami with regions}
    \label{tab:SCBsami}
    \begin{tabular}{L{3cm}|L{6.5cm}}
      \textbf{\gls{ISO} 639-2 Code} &\textbf{Swedish name of Language}\\
      \hline
SMI &   Samiska\\
SIA &   Samiska, Akkalasamiska\\
SMN &   Samiska, Enaresamiska\\
SJK &   Samiska, Kemisamiska\\
SJD &   Samiska, Kildinsamiska\\
SMJ &   Samiska, Lulesamiska(Luleå)\\
SJE &   Samiska, Pitesamiska\\
SMS &   Samiska, Skoltsamiska\\
SMA &  Samiska, sydsamiska (södra)\\
SJT &  Samiska, Tersamiska\\
SJU &  Umesamiska\\
\end{tabular}
\end{table}
\clearpage
\begin{table}[!htp]
    \caption{Languages already implemented}
    \label{tab:languagesImplemented}
    \begin{footnotesize}

    \begin{tabular}{L{3cm}|L{3cm}|L{6.5cm}}
      \textbf{\gls{ISO} 639-2 Code} & \textbf{\gls{ISO} 639-1 Code} & \textbf{English name of Language}\\
      \hline
alb & sq & Albanian\\
ara & ar & Arabic\\
bel & be & Belarusian\\
bul & bg & Bulgarian\\
cat & ca & Catalan; Valencian\\
chi & zh & Chinese\\
cze & cs & Czech\\
dan & da & Danish\\
dut & nl & Dutch; Flemish\\
eng & en & English\\
est & et & Estonian\\
fin & fi & Finnish\\
fre & fr & French\\
ger & de & German\\
gre & el & Greek, Modern (1453-)\\
heb & he & Hebrew\\
hin & hi & Hindi\\
hun & hu & Hungarian\\
ice & is & Icelandic\\
ita & it & Italian\\
jpn & ja & Japanese\\
lat & la & Latin\\
lav & lv & Latvian\\
lit & lt & Lithuanian\\
nno & nn & Norwegian Nynorsk; Nynorsk, Norwegian\\
nob & nb & Bokmål, Norwegian; Norwegian Bokmål\\
nor & no & Norwegian\\
per & fa & Persian\\
pol & pl & Polish\\
por & pt & Portuguese\\
rum & ro & Romanian; Moldavian; Moldovan\\
rus & ru & Russian\\
sme & se & Northern Sami\\
srp & sr & Serbian\\
slo & sk & Slovak\\
spa & es & Spanish; Castilian\\
swe & sv & Swedish\\
swa & sw & Swahili\\
tur & tr & Turkish\\
ukr & uk & Ukrainian\\
vie & vi & Vietnamese\\
yid & yi & Yiddish\\
\hline
   \end{tabular}  
    \end{footnotesize}
\end{table}
\FloatBarrier

\begin{table}[!ht]
    \caption{Missing \gls{ISO} 639-1 Code from the Library of Congress list}
    \label{tab:languagesMissing}
    \begin{tabular}{L{3cm}|L{3cm}|L{6.5cm}}
      \textbf{\gls{ISO} 639-2 Code} & \textbf{\gls{ISO} 639-1 Code} & \textbf{English name of Language}\\
      \hline
fit & ??? & Meänkieli\\
\hline
   \end{tabular}
\end{table}
\FloatBarrier

As to the language tag to use for Meänkieli, one can simply skip the \gls{ISO} 639-1 Code and use "fit" - as done in \url{https://www.msb.se/sv/om-msb/information-pa-andra-sprak/meankieli/om-msb/}.
However, it seems strange that there is no \gls{ISO} 639-1 Code. Interestingly, they use 'smi' for Lule Sami (\foreignlanguage{swedish}{lulesamiska}) in \url{https://www.msb.se/sv/om-msb/information-pa-andra-sprak/lulesamiska/om-msb/}.
\generalExpl{I have added a skeleton abstract for \foreignlanguage{swedish}{Meänkieli}. I had help from a former professor at Stockholm's university and sprakcentrummeankieli@isof.se for the abstract and keywords in Meänkieli (although there is a small difference between their spellings of two words).}

\begin{longtable}{|L{3cm}|L{3cm}|L{5cm}|}
\caption{Remaining languages to implement}
\label{tab:languagesNotImplemented}\\
\hline \multicolumn{1}{|c|}{\textbf{\gls{ISO} 639-2 Code}} & \multicolumn{1}{c|}{\textbf{\gls{ISO} 639-1 Code}} & \multicolumn{1}{c|}{\textbf{English name of Language}} \\
\endfirsthead
\multicolumn{3}{c}%
{{\bfseries \tablename\ \thetable{} -- continued from previous page}} \\
\hline \multicolumn{1}{|c|}{\textbf{\gls{ISO} 639-2 Code}} & \multicolumn{1}{c|}{\textbf{\gls{ISO} 639-1 Code}} & \multicolumn{1}{c|}{\textbf{English name of Language}} \\
\hline 
\endhead
\hline \multicolumn{3}{|r|}{{Continued on next page}} \\
\hline
\endfoot
\hline
\hline
\endlastfoot
aar & aa & Afar\\
abk & ab & Abkhazian\\
afr & af & Afrikaans\\
aka & ak & Akan\\
amh & am & Amharic\\
arg & an & Aragonese\\
arm & hy & Armenian\\
asm & as & Assamese\\
ava & av & Avaric\\
ave & ae & Avestan\\
aym & ay & Aymara\\
aze & az & Azerbaijani\\
bak & ba & Bashkir\\
bam & bm & Bambara\\
baq & eu & Basque\\
ben & bn & Bengali\\
bih & bh & Bihari languages\\
bis & bi & Bislama\\
bos & bs & Bosnian\\
bre & br & Breton\\
bur & my & Burmese\\
cha & ch & Chamorro\\
che & ce & Chechen\\
chu & cu & Church Slavic; Old Slavonic; Church Slavonic; Old Bulgarian; Old Church Slavonic\\
chv & cv & Chuvash\\
cor & kw & Cornish\\
cos & co & Corsican\\
cre & cr & Cree\\
div & dv & Divehi; Dhivehi; Maldivian\\
dzo & dz & Dzongkha\\
epo & eo & Esperanto\\
ewe & ee & Ewe\\
fao & fo & Faroese\\
fij & fj & Fijian\\
fry & fy & Western Frisian\\
ful & ff & Fulah\\
geo & ka & Georgian\\
gla & gd & Gaelic; Scottish\\
gle & ga & Irish\\
glg & gl & Galician\\
glv & gv & Manx\\
grn & gn & Guarani\\
guj & gu & Gujarati\\
hat & ht & Haitian; Haitian\\
hau & ha & Hausa\\
her & hz & Herero\\
hmo & ho & Hiri Motu\\
hrv & hr & Croatian\\
ibo & ig & Igbo\\
ido & io & Ido\\
iii & ii & Sichuan Yi; Nuosu\\
iku & iu & Inuktitut\\
ile & ie & Interlingue; Occidental\\
ina & ia & Interlingua (International Auxiliary Language Association)\\
ind & id & Indonesian\\
ipk & ik & Inupiaq\\
jav & jv & Javanese\\
kal & kl & Kalaallisut; Greenlandic\\
kan & kn & Kannada\\
kas & ks & Kashmiri\\
kau & kr & Kanuri\\
kaz & kk & Kazakh\\
khm & km & Central Khmer\\
kik & ki & Kikuyu; Gikuyu\\
kin & rw & Kinyarwanda\\
kir & ky & Kirghiz; Kyrgyz\\
kom & kv & Komi\\
kon & kg & Kongo\\
kor & ko & Korean\\
kua & kj & Kuanyama; Kwanyama\\
kur & ku & Kurdish\\
lao & lo & Lao\\
lim & li & Limburgan; Limburger; Limburgish\\
lin & ln & Lingala\\
ltz & lb & Luxembourgish; Letzeburgesch\\
lub & lu & Luba-Katanga\\
lug & lg & Ganda\\
mac & mk & Macedonian\\
mah & mh & Marshallese\\
mal & ml & Malayalam\\
mao & mi & Maori\\
mar & mr & Marathi\\
may & ms & Malay\\
mlg & mg & Malagasy\\
mlt & mt & Maltese\\
mon & mn & Mongolian\\
nau & na & Nauru\\
nav & nv & Navajo; Navaho\\
nbl & nr & Ndebele, South; South Ndebele\\
nde & nd & Ndebele, North; North Ndebele\\
ndo & ng & Ndonga\\
nep & ne & Nepali\\
nya & ny & Chichewa; Chewa; Nyanja\\
oci & oc & Occitan (post 1500)\\
oji & oj & Ojibwa\\
ori & or & Oriya\\
orm & om & Oromo\\
oss & os & Ossetian; Ossetic\\
pan & pa & Panjabi; Punjabi\\
pli & pi & Pali\\
pus & ps & Pushto; Pashto\\
que & qu & Quechua\\
roh & rm & Romansh\\
run & rn & Rundi\\
sag & sg & Sango\\
san & sa & Sanskrit\\
sin & si & Sinhala; Sinhalese\\
slv & sl & Slovenian\\
smo & sm & Samoan\\
sna & sn & Shona\\
snd & sd & Sindhi\\
som & so & Somali\\
sot & st & Sotho, Southern\\
srd & sc & Sardinian\\
ssw & ss & Swati\\
sun & su & Sundanese\\
tah & ty & Tahitian\\
tam & ta & Tamil\\
tat & tt & Tatar\\
tel & te & Telugu\\
tgk & tg & Tajik\\
tgl & tl & Tagalog\\
tha & th & Thai\\
tib & bo & Tibetan\\
tir & ti & Tigrinya\\
ton & to & Tonga (Tonga Islands)\\
tsn & tn & Tswana\\
tso & ts & Tsonga\\
tuk & tk & Turkmen\\
twi & tw & Twi\\
uig & ug & Uighur; Uyghur\\
urd & ur & Urdu\\
uzb & uz & Uzbek\\
ven & ve & Venda\\
vol & vo & Volapük\\
wel & cy & Welsh\\
wln & wa & Walloon\\
wol & wo & Wolof\\
xho & xh & Xhosa\\
yor & yo & Yoruba\\
zha & za & Zhuang; Chuang\\
zul & zu & Zulu\\
\end{longtable}
\FloatBarrier
\clearpage

\section{Miscellaneous Unicode symbols}
\label{sec:unicodeSymbols}
A great reference that helped identify the relevant fonts was \url{https://codepoints.net/}; see, for example, \url{https://codepoints.net/U+26FF}.
\noindent
\begingroup
\captionsetup{type=listing,hypcap=false}
\caption{U+2600--U+26FF Miscellaneous Symbols}
\label{lst:unicodeSymbols}
\vspace{\abovecaptionskip}
\begin{minted}[breaklines, escapeinside=||]{tex}
%%%%%%%%%%%%%%%%%%%%%%%%%%%%%%%%%%%%%%%%%%%%%%%%%%%%%%%%%%%%%%%%%%%%%%
% Miscellaneous Symbols, U+2600 - U+26FF
% The characters below had been used in a title, subtitle, abstract, or keywords on or before 2025-05-28
\DeclareUnicodeFallback{2600}{\mystixmathfont}{^^^^2600} % Black Sun with Rays - ☀
\DeclareUnicodeFallback{2601}{\mystixmathfont}{^^^^2601} % Cloud - ☁
\DeclareUnicodeFallback{2602}{\mystixmathfont}{^^^^2602} % Umbrella - ☂
\DeclareUnicodeFallback{2603}{\mystixmathfont}{^^^^2603} % Snowman - ☃
\DeclareUnicodeFallback{2604}{\mystixmathfont}{^^^^2604} % Comet - ☄
\DeclareUnicodeFallback{2605}{\mystixmathfont}{^^^^2605} % Black Star - ★
\DeclareUnicodeFallback{2606}{\mystixmathfont}{^^^^2606} % White Star - ☆
\DeclareUnicodeFallback{2607}{\mystixmathfont}{^^^^2607} % Lightning - ☇
\DeclareUnicodeFallback{2608}{\mystixmathfont}{^^^^2608} % Thunderstorm - ☈
\DeclareUnicodeFallback{2609}{\mystixmathfont}{^^^^2609} % Sun - ☉
\DeclareUnicodeFallback{260A}{\mystixmathfont}{^^^^260a} % Ascending Node - ☊
\DeclareUnicodeFallback{260B}{\mystixmathfont}{^^^^260b} % Descending Node - ☋
\DeclareUnicodeFallback{260C}{\mystixmathfont}{^^^^260c} % Conjunction - ☌
\DeclareUnicodeFallback{260D}{\mystixmathfont}{^^^^260d} % Opposition - ☍
\DeclareUnicodeFallback{260E}{\mystixmathfont}{^^^^260e} % Black Telephone - ☎
\DeclareUnicodeFallback{260F}{\mystixmathfont}{^^^^260f} % White Telephone - ☏

\end{minted}
\newpage
\begin{minted}[breaklines, escapeinside=||]{tex}
\DeclareUnicodeFallback{2610}{\mystixmathfont}{^^^^2610} % Ballot Box - ☐
\DeclareUnicodeFallback{2611}{\mystixmathfont}{^^^^2611} % Ballot Box with Check - ☑
\DeclareUnicodeFallback{2612}{\mystixmathfont}{^^^^2612} % Ballot Box with X - ☒
\DeclareUnicodeFallback{2613}{\mystixmathfont}{^^^^2613} % Saltire - ☓
\DeclareUnicodeFallback{2614}{\NotoSansSymbolsTwoFont}{^^^^2614} % Umbrella with Rain Drops - ☔
\DeclareUnicodeFallback{2615}{\NotoEmojiFont}{^^^^2615} % Hot Beverage - ☕
\DeclareUnicodeFallback{2616}{\NotoSansJPFont}{^^^^2616} % White Shogi Piece - ☖
\DeclareUnicodeFallback{2617}{\NotoSansJPFont}{^^^^2617} % Black Shogi Piece - ☗
\DeclareUnicodeFallback{2618}{\NotoEmojiFont}{^^^^2618} % Shamrock - ☘
\DeclareUnicodeFallback{2619}{\NotoSansSymbolsTwoFont}{^^^^2619} % Reversed Rotated Floral Heart Bullet - ☙
\DeclareUnicodeFallback{261A}{\mystixmathfont}{^^^^261a} % Black Left Pointing Index - ☚
\DeclareUnicodeFallback{261B}{\mystixmathfont}{^^^^261b} % Black Right Pointing Index - ☛
\DeclareUnicodeFallback{261C}{\mystixmathfont}{^^^^261c} % White Left Pointing Index - ☜
\DeclareUnicodeFallback{261D}{\mystixmathfont}{^^^^261d} % White Up Pointing Index - ☝
\DeclareUnicodeFallback{261E}{\mystixmathfont}{^^^^261e} % White Right Pointing Index - ☞
\DeclareUnicodeFallback{261F}{\mystixmathfont}{^^^^261f} % White Down Pointing Index - ☟

\end{minted}
\newpage
\begin{minted}[breaklines, escapeinside=||]{tex}
\DeclareUnicodeFallback{2620}{\mystixmathfont}{^^^^2620} % Skull and Crossbones - ☠
\DeclareUnicodeFallback{2621}{\mystixmathfont}{^^^^2621} % Caution Sign - ☡
\DeclareUnicodeFallback{2622}{\mystixmathfont}{^^^^2622} % Radioactive Sign - ☢
\DeclareUnicodeFallback{2623}{\mystixmathfont}{^^^^2623} % Biohazard Sign - ☣
\DeclareUnicodeFallback{2624}{\mystixmathfont}{^^^^2624} % Caduceus - ☤
\DeclareUnicodeFallback{2625}{\mystixmathfont}{^^^^2625} % Ankh - ☥
\DeclareUnicodeFallback{2626}{\mystixmathfont}{^^^^2626} % Orthodox Cross - ☦
\DeclareUnicodeFallback{2627}{\mystixmathfont}{^^^^2627} % Chi Rho - ☧
\DeclareUnicodeFallback{2628}{\mystixmathfont}{^^^^2628} % Cross of Lorraine - ☨
\DeclareUnicodeFallback{2629}{\mystixmathfont}{^^^^2629} % Cross of Jerusalem - ☩
\DeclareUnicodeFallback{262A}{\mystixmathfont}{^^^^262a} % Star and Crescent - ☪
\DeclareUnicodeFallback{262B}{\mystixmathfont}{^^^^262b} % Farsi Symbol - ☫
\DeclareUnicodeFallback{262C}{\mystixmathfont}{^^^^262c} % Adi Shakti - ☬
\DeclareUnicodeFallback{262D}{\mystixmathfont}{^^^^262d} % Hammer and Sickle - ☭
\DeclareUnicodeFallback{262E}{\mystixmathfont}{^^^^262e} % Peace Symbol - ☮
\DeclareUnicodeFallback{262F}{\mystixmathfont}{^^^^262f} % Yin Yang - ☯

\end{minted}
\newpage
\begin{minted}[breaklines, escapeinside=||]{tex}
\DeclareUnicodeFallback{2630}{\mystixmathfont}{^^^^2630} % Trigram For Heaven - ☰
\DeclareUnicodeFallback{2631}{\mystixmathfont}{^^^^2631} % Trigram For Lake - ☱
\DeclareUnicodeFallback{2632}{\mystixmathfont}{^^^^2632} % Trigram For Fire - ☲
\DeclareUnicodeFallback{2633}{\mystixmathfont}{^^^^2633} % Trigram For Thunder - ☳
\DeclareUnicodeFallback{2634}{\mystixmathfont}{^^^^2634} % Trigram For Wind - ☴
\DeclareUnicodeFallback{2635}{\mystixmathfont}{^^^^2635} % Trigram For Water- ☵
\DeclareUnicodeFallback{2636}{\mystixmathfont}{^^^^2636} % Trigram For Mountain - ☶
\DeclareUnicodeFallback{2637}{\mystixmathfont}{^^^^2637} % Trigram For Earth - ☷
\DeclareUnicodeFallback{2638}{\mystixmathfont}{^^^^2638} % Wheel of Dharma - ☸
\DeclareUnicodeFallback{2639}{\mystixmathfont}{^^^^2639} % White Frowning Face - ☹
\DeclareUnicodeFallback{263A}{\NotoEmojiFont}{^^^^263a} % White Smiling Face - ☺
\DeclareUnicodeFallback{263B}{\mystixmathfont}{^^^^263b} % Black Smiling Face - ☻
\DeclareUnicodeFallback{263C}{\mystixmathfont}{^^^^263c} % White Sun with Rays - ☼
\DeclareUnicodeFallback{263D}{\mystixmathfont}{^^^^263d} % First Quarter Moon- ☽
\DeclareUnicodeFallback{263E}{\mystixmathfont}{^^^^263e} % PLast Quarter Moon - ☾
\DeclareUnicodeFallback{263F}{\mystixmathfont}{^^^^263f} % Mercur - ☿

\end{minted}
\newpage
\begin{minted}[breaklines, escapeinside=||]{tex}
\DeclareUnicodeFallback{2640}{\mystixmathfont}{^^^^2640} % female symbol - ♀
\DeclareUnicodeFallback{2641}{\mystixmathfont}{^^^^2641} % Earth - ♁
\DeclareUnicodeFallback{2642}{\mystixmathfont}{^^^^2642} % Male symbol - ♂
\DeclareUnicodeFallback{2643}{\mystixmathfont}{^^^^2643} % Jupiter - ♃
\DeclareUnicodeFallback{2644}{\mystixmathfont}{^^^^2644} % Saturn - ♄
\DeclareUnicodeFallback{2645}{\mystixmathfont}{^^^^2645} % Uranus- ♅
\DeclareUnicodeFallback{2646}{\mystixmathfont}{^^^^2646} % Neptune - ♆
\DeclareUnicodeFallback{2647}{\mystixmathfont}{^^^^2647} % Pluto - ♇
\DeclareUnicodeFallback{2648}{\mystixmathfont}{^^^^2648} % Aries - ♈
\DeclareUnicodeFallback{2649}{\mystixmathfont}{^^^^2649} % Taurus - ♉
\DeclareUnicodeFallback{264A}{\mystixmathfont}{^^^^264a} % Gemini - ♊
\DeclareUnicodeFallback{264B}{\mystixmathfont}{^^^^264b} % Cancer - ♋
\DeclareUnicodeFallback{264C}{\mystixmathfont}{^^^^264c} % Leo - ♌
\DeclareUnicodeFallback{264D}{\mystixmathfont}{^^^^264d} % Virgo - ♍
\DeclareUnicodeFallback{264E}{\mystixmathfont}{^^^^264e} % Libra - ♎
\DeclareUnicodeFallback{264F}{\mystixmathfont}{^^^^264f} % Scorpiu - ♏

\end{minted}
\newpage
\begin{minted}[breaklines, escapeinside=||]{tex}
\DeclareUnicodeFallback{2650}{\mystixmathfont}{^^^^2650} % Sagittarius - ♐
\DeclareUnicodeFallback{2651}{\mystixmathfont}{^^^^2651} % Capricorn - ♑
\DeclareUnicodeFallback{2652}{\mystixmathfont}{^^^^2652} % Aquarius - ♒
\DeclareUnicodeFallback{2653}{\mystixmathfont}{^^^^2653} % Pisces - ♓
\DeclareUnicodeFallback{2654}{\mystixmathfont}{^^^^2654} % White Chess King - ♔
\DeclareUnicodeFallback{2655}{\mystixmathfont}{^^^^2655} % White Chess Queen- ♕
\DeclareUnicodeFallback{2656}{\mystixmathfont}{^^^^2656} % White Chess Rook - ♖
\DeclareUnicodeFallback{2657}{\mystixmathfont}{^^^^2657} % White Chess Bishop - ♗
\DeclareUnicodeFallback{2658}{\mystixmathfont}{^^^^2658} % White Chess Knight - ♘
\DeclareUnicodeFallback{2659}{\mystixmathfont}{^^^^2659} % White Chess Pawn - ♙
\DeclareUnicodeFallback{265A}{\mystixmathfont}{^^^^265a} % Black Chess King - ♚
\DeclareUnicodeFallback{265B}{\mystixmathfont}{^^^^265b} % Black Chess Queen - ♛
\DeclareUnicodeFallback{265C}{\mystixmathfont}{^^^^265c} % Black Chess Rook - ♜
\DeclareUnicodeFallback{265D}{\mystixmathfont}{^^^^265d} % Black Chess Bishop - ♝
\DeclareUnicodeFallback{265E}{\mystixmathfont}{^^^^265e} % Black Chess Knight - ♞
\DeclareUnicodeFallback{265F}{\mystixmathfont}{^^^^265f} % Black Chess Pawn - ♟

\end{minted}
\newpage
\begin{minted}[breaklines, escapeinside=||]{tex}
\DeclareUnicodeFallback{2660}{\mystixmathfont}{^^^^2660} % Black Spade Suit - ♠
\DeclareUnicodeFallback{2661}{\mystixmathfont}{^^^^2661} % White Heart Suit - ♡
\DeclareUnicodeFallback{2662}{\mystixmathfont}{^^^^2662} % White Diamond Suit - ♢
\DeclareUnicodeFallback{2663}{\mystixmathfont}{^^^^2663} % Black Club Suit - ♣
\DeclareUnicodeFallback{2664}{\mystixmathfont}{^^^^2664} % White Spade Suit - ♤
\DeclareUnicodeFallback{2665}{\mystixmathfont}{^^^^2665} % Black Heart Suit- ♥
\DeclareUnicodeFallback{2666}{\mystixmathfont}{^^^^2666} % Black Diamond Suit - ♦
\DeclareUnicodeFallback{2667}{\mystixmathfont}{^^^^2667} % White Club Suit - ♧
\DeclareUnicodeFallback{2668}{\mystixmathfont}{^^^^2668} % Hot Springs - ♨
\DeclareUnicodeFallback{2669}{\mystixmathfont}{^^^^2669} % Quarter Note - ♩
\DeclareUnicodeFallback{266A}{\mystixmathfont}{^^^^266a} % Eighth Note - ♪
\DeclareUnicodeFallback{266B}{\mystixmathfont}{^^^^266b} % Beamed Eighth Notes - ♫
\DeclareUnicodeFallback{266C}{\mystixmathfont}{^^^^266c} % Beamed Sixteenth Notes - ♬
\DeclareUnicodeFallback{266D}{\mystixmathfont}{^^^^266d} % Music Flat Sign - ♭
\DeclareUnicodeFallback{266E}{\mystixmathfont}{^^^^266e} % Music Natural Sign - ♮
\DeclareUnicodeFallback{266F}{\mystixmathfont}{^^^^266f} % Music Sharp Sign - ♯

\end{minted}
\newpage
\begin{minted}[breaklines, escapeinside=||]{tex}
\DeclareUnicodeFallback{2670}{\mystixmathfont}{^^^^2670} % West Syriac Cross - ♰
\DeclareUnicodeFallback{2671}{\mystixmathfont}{^^^^2671} % East Syriac Cross - ♱
\DeclareUnicodeFallback{2672}{\NotoSansSymbolsFont}{^^^^2672} % Universal Recycling Symbol - ♲
\DeclareUnicodeFallback{2673}{\NotoSansSymbolsFont}{^^^^2673} % Recycling Symbol For Type-1 Plastics - ♳
\DeclareUnicodeFallback{2674}{\NotoSansSymbolsFont}{^^^^2674} % Recycling Symbol For Type-2 Plastics - ♴
\DeclareUnicodeFallback{2675}{\NotoSansSymbolsFont}{^^^^2675} % Recycling Symbol For Type-3 Plastics- ♵
\DeclareUnicodeFallback{2676}{\NotoSansSymbolsFont}{^^^^2676} % Recycling Symbol For Type-4 Plastics - ♶
\DeclareUnicodeFallback{2677}{\NotoSansSymbolsFont}{^^^^2677} % Recycling Symbol For Type-5 Plastics - ♷
\DeclareUnicodeFallback{2678}{\NotoSansSymbolsFont}{^^^^2678} % Recycling Symbol For Type-6 Plastics - ♸
\DeclareUnicodeFallback{2679}{\NotoSansSymbolsFont}{^^^^2679} % Recycling Symbol For Type-7 Plastics - ♹
\DeclareUnicodeFallback{267A}{\mystixmathfont}{^^^^267a} % Recycling Symbol For Generic Materials - ♺
\DeclareUnicodeFallback{267B}{\mystixmathfont}{^^^^267b} % Black Universal Recycling Symbol - ♻
\DeclareUnicodeFallback{267C}{\NotoSansSymbolsFont}{^^^^267c} % Recycled Paper Symbol - ♼
\DeclareUnicodeFallback{267D}{\NotoSansSymbolsFont}{^^^^267d} % Partially-Recycled Paper Symbol - ♽
\DeclareUnicodeFallback{267E}{\mystixmathfont}{^^^^267e} % Permanent Paper Sign - ♾
\DeclareUnicodeFallback{267F}{\NotoEmojiFont}{^^^^267f} % Wheelchair Symbol - ♿

\end{minted}
\newpage
\begin{minted}[breaklines, escapeinside=||]{tex}
\DeclareUnicodeFallback{2680}{\mystixmathfont}{^^^^2680} % Die Face-1 - ⚀
\DeclareUnicodeFallback{2681}{\mystixmathfont}{^^^^2681} % Die Face-2 - ⚁
\DeclareUnicodeFallback{2682}{\mystixmathfont}{^^^^2682} % Die Face-3 - ⚂
\DeclareUnicodeFallback{2683}{\mystixmathfont}{^^^^2683} % Die Face-4 - ⚃
\DeclareUnicodeFallback{2684}{\mystixmathfont}{^^^^2684} % Die Face-5 - ⚄
\DeclareUnicodeFallback{2685}{\mystixmathfont}{^^^^2685} % Die Face-6- ⚅
\DeclareUnicodeFallback{2686}{\mystixmathfont}{^^^^2686} % White Circle with Dot Right - ⚆
\DeclareUnicodeFallback{2687}{\mystixmathfont}{^^^^2687} % White Circle with Two Dots - ⚇
\DeclareUnicodeFallback{2688}{\mystixmathfont}{^^^^2688} % Black Circle with White Dot Right - ⚈
\DeclareUnicodeFallback{2689}{\mystixmathfont}{^^^^2689} % Black Circle with Two White Dots - ⚉
\DeclareUnicodeFallback{268A}{\mystixmathfont}{^^^^268a} % Monogram For Yang - ⚊
\DeclareUnicodeFallback{268B}{\mystixmathfont}{^^^^268b} % Monogram For Yin - ⚋
\DeclareUnicodeFallback{268C}{\mystixmathfont}{^^^^268c} % Digram For Greater Yang - ⚌
\DeclareUnicodeFallback{268D}{\mystixmathfont}{^^^^268d} % Digram For Lesser Yin - ⚍
\DeclareUnicodeFallback{268E}{\mystixmathfont}{^^^^268e} % Digram For Lesser Yang - ⚎
\DeclareUnicodeFallback{268F}{\mystixmathfont}{^^^^268f} % Digram For Greater Yin - ⚏

\end{minted}
\newpage
\begin{minted}[breaklines, escapeinside=||]{tex}
\DeclareUnicodeFallback{2690}{\mystixmathfont}{^^^^2690} % White Flag - ⚐
\DeclareUnicodeFallback{2691}{\mystixmathfont}{^^^^2691} % Black Flag - ⚑
\DeclareUnicodeFallback{2692}{\NotoSansSymbolsFont}{^^^^2692} % Hammer and Pick - ⚒
\DeclareUnicodeFallback{2693}{\NotoSansSymbolsFont}{^^^^2693} % Anchor - ⚓
\DeclareUnicodeFallback{2694}{\NotoSansSymbolsFont}{^^^^2694} % Crossed Swords - ⚔
\DeclareUnicodeFallback{2695}{\NotoSansSymbolsFont}{^^^^2695} % Staff of Aesculapius - ⚕
\DeclareUnicodeFallback{2696}{\NotoSansSymbolsFont}{^^^^2696} % Scales - ⚖
\DeclareUnicodeFallback{2697}{\NotoSansSymbolsFont}{^^^^2697} % Alembic - ⚗
\DeclareUnicodeFallback{2698}{\NotoSansSymbolsFont}{^^^^2698} % Flower - ⚘
\DeclareUnicodeFallback{2699}{\NotoSansSymbolsFont}{^^^^2699} % Gear - ⚙
\DeclareUnicodeFallback{269A}{\mystixmathfont}{^^^^269a} % Staff of Hermes - ⚚
\DeclareUnicodeFallback{269B}{\mystixmathfont}{^^^^269b} % Atom Symbol - ⚛
\DeclareUnicodeFallback{269C}{\mystixmathfont}{^^^^269c} % Fleur-De-Lis - ⚜
\DeclareUnicodeFallback{269D}{\mystixmathfont}{^^^^269d} % Outlined White Star - ⚝
\DeclareUnicodeFallback{269E}{\NotoSansSymbolsTwoFont}{^^^^269e} % Three Lines Converging Right - ⚞
\DeclareUnicodeFallback{269F}{\NotoSansSymbolsTwoFont}{^^^^269f} % Three Lines Converging Left - ⚟

\end{minted}
\newpage
\begin{minted}[breaklines, escapeinside=||]{tex}
\DeclareUnicodeFallback{26A0}{\mystixmathfont}{^^^^26a0} % Warning Sign - ⚠
\DeclareUnicodeFallback{26A1}{\mystixmathfont}{^^^^26a1} % High Voltage Sign - ⚡
\DeclareUnicodeFallback{26A2}{\mystixmathfont}{^^^^26a2} % Doubled Female Sign - ⚢
\DeclareUnicodeFallback{26A3}{\mystixmathfont}{^^^^26a3} % Doubled Male Sign - ⚣
\DeclareUnicodeFallback{26A4}{\mystixmathfont}{^^^^26a4} % Interlocked Female and Male Sign - ⚤
\DeclareUnicodeFallback{26A5}{\mystixmathfont}{^^^^26a5} % Male and Female Sign - ⚥
\DeclareUnicodeFallback{26A6}{\mystixmathfont}{^^^^26a6} % Male with Stroke Sign - ⚦
\DeclareUnicodeFallback{26A7}{\mystixmathfont}{^^^^26a7} % Male with Stroke and Male and Female Sign - ⚧
\DeclareUnicodeFallback{26A8}{\mystixmathfont}{^^^^26a8} % Vertical Male with Stroke Sign - ⚨
\DeclareUnicodeFallback{26A9}{\mystixmathfont}{^^^^26a9} % Horizontal Male with Stroke Sign - ⚩
\DeclareUnicodeFallback{26AA}{\mystixmathfont}{^^^^26aa} % Medium White Circle - ⚪
\DeclareUnicodeFallback{26AB}{\mystixmathfont}{^^^^26ab} % Medium Black Circle - ⚫
\DeclareUnicodeFallback{26AC}{\mystixmathfont}{^^^^26ac} % Medium Small White Circle - ⚬
\DeclareUnicodeFallback{26AD}{\mystixmathfont}{^^^^26ad} % Marriage Symbol - ⚭
\DeclareUnicodeFallback{26AE}{\mystixmathfont}{^^^^26ae} % Divorce Symbol - ⚮
\DeclareUnicodeFallback{26AF}{\mystixmathfont}{^^^^26af} % Unmarried Partnership Symbol - ⚯

\end{minted}
\newpage
\begin{minted}[breaklines, escapeinside=||]{tex}
\DeclareUnicodeFallback{26B0}{\NotoEmojiFont}{^^^^26b0} % Coffin - ⚰
\DeclareUnicodeFallback{26B1}{\NotoEmojiFont}{^^^^26b1} % Funeral Urn - ⚱
\DeclareUnicodeFallback{26B2}{\mystixmathfont}{^^^^26b2} % Neuter - ⚲
\DeclareUnicodeFallback{26B3}{\mystixmathfont}{^^^^26b3} % Ceres - ⚳
\DeclareUnicodeFallback{26B4}{\mystixmathfont}{^^^^26b4} % Pallas - ⚴
\DeclareUnicodeFallback{26B5}{\mystixmathfont}{^^^^26b5} % Juno - ⚵
\DeclareUnicodeFallback{26B6}{\mystixmathfont}{^^^^26b6} % Vesta - ⚶
\DeclareUnicodeFallback{26B7}{\mystixmathfont}{^^^^26b7} % Chiron - ⚷
\DeclareUnicodeFallback{26B8}{\mystixmathfont}{^^^^26b8} % Black Moon Lilith - ⚸
\DeclareUnicodeFallback{26B9}{\mystixmathfont}{^^^^26b9} % Sextile - ⚹
\DeclareUnicodeFallback{26BA}{\mystixmathfont}{^^^^26ba} % Semisextile - ⚺
\DeclareUnicodeFallback{26BB}{\mystixmathfont}{^^^^26bb} % Quincunx - ⚻
\DeclareUnicodeFallback{26BC}{\mystixmathfont}{^^^^26bc} % Sesquiquadrate - ⚼
\DeclareUnicodeFallback{26BD}{\NotoEmojiFont}{^^^^26bd} % Soccer Ball - ⚽
\DeclareUnicodeFallback{26BE}{\NotoEmojiFont}{^^^^26be} % Baseball - ⚾
\DeclareUnicodeFallback{26BF}{\NotoSansSymbolsTwoFont}{^^^^26bf} % Squared Key - ⚿

\end{minted}
\newpage
\begin{minted}[breaklines, escapeinside=||]{tex}
\DeclareUnicodeFallback{26C0}{\NotoSansSymbolsTwoFont}{^^^^26c0} % White Draughts Man - ⛀
\DeclareUnicodeFallback{26C1}{\NotoSansSymbolsTwoFont}{^^^^26c1} % White Draughts King - ⛁
\DeclareUnicodeFallback{26C2}{\NotoSansSymbolsTwoFont}{^^^^26c2} % Black Draughts Man - ⛂
\DeclareUnicodeFallback{26C3}{\NotoSansSymbolsTwoFont}{^^^^26c3} % Black Draughts King - ⛃
\DeclareUnicodeFallback{26C4}{\NotoEmojiFont}{^^^^26c4} % Snowman Without Snow - ⛄
\DeclareUnicodeFallback{26C5}{\NotoEmojiFont}{^^^^26c5} % Sun Behind Cloud - ⛅
\DeclareUnicodeFallback{26C6}{\NotoSansSymbolsTwoFont}{^^^^26c6} % Rain - ⛆
\DeclareUnicodeFallback{26C7}{\NotoSansSymbolsTwoFont}{^^^^26c7} % Black Snowman - ⛇
\DeclareUnicodeFallback{26C8}{\NotoEmojiFont}{^^^^26c8} % Thunder Cloud and Rain - ⛈
\DeclareUnicodeFallback{26C9}{\NotoSansSymbolsTwoFont}{^^^^26c9} % Turned White Shogi Piece - ⛉
\DeclareUnicodeFallback{26CA}{\NotoSansSymbolsTwoFont}{^^^^26ca} % Turned Black Shogi Piece - ⛊
\DeclareUnicodeFallback{26CB}{\mystixmathfont}{^^^^26cb} % White Diamond In Square - ⛋
\DeclareUnicodeFallback{26CC}{\mystixmathfont}{^^^^26cc} % Crossing Lanes - ⛌
\DeclareUnicodeFallback{26CD}{\NotoSansSymbolsTwoFont}{^^^^26cd} % Disabled Car - ⛍
\DeclareUnicodeFallback{26CE}{\NotoEmojiFont}{^^^^26ce} % Ophiuchus - ⛎
\DeclareUnicodeFallback{26CF}{\NotoSansSymbolsTwoFont}{^^^^26cf} % Pick - ⛏

\end{minted}
\newpage
\begin{minted}[breaklines, escapeinside=||]{tex}
\DeclareUnicodeFallback{26D0}{\NotoSansSymbolsTwoFont}{^^^^26d0} % Car Sliding - ⛐
\DeclareUnicodeFallback{26D1}{\NotoSansSymbolsTwoFont}{^^^^26d1} % Helmet with White Cross - ⛑
\DeclareUnicodeFallback{26D2}{\NotoSansSymbolsTwoFont}{^^^^26d2} % Circled Crossing Lanes - ⛒
\DeclareUnicodeFallback{26D3}{\NotoSansSymbolsTwoFont}{^^^^26d3} % Chains - ⛓
\DeclareUnicodeFallback{26D4}{\NotoSansSymbolsTwoFont}{^^^^26d4} % No Entry - ⛔
\DeclareUnicodeFallback{26D5}{\NotoSansSymbolsTwoFont}{^^^^26d5} % Alternate One-Way Left Way Traffic - ⛕
\DeclareUnicodeFallback{26D6}{\NotoSansSymbolsTwoFont}{^^^^26d6} % Black Two-Way Left Way Traffic - ⛖
\DeclareUnicodeFallback{26D7}{\NotoSansSymbolsTwoFont}{^^^^26d7} % White Two-Way Left Way Traffic - ⛗
\DeclareUnicodeFallback{26D8}{\NotoSansSymbolsTwoFont}{^^^^26d8} % Black Left Lane Merge - ⛘
\DeclareUnicodeFallback{26D9}{\NotoSansSymbolsTwoFont}{^^^^26d9} % White Left Lane Merge - ⛙
\DeclareUnicodeFallback{26DA}{\NotoSansSymbolsTwoFont}{^^^^26da} % Drive Slow Sign - ⛚
\DeclareUnicodeFallback{26DB}{\NotoSansSymbolsTwoFont}{^^^^26db} % Heavy White Down-Pointing Triangle - ⛛
\DeclareUnicodeFallback{26DC}{\NotoSansSymbolsTwoFont}{^^^^26dc} % Left Closed Entry - ⛜
\DeclareUnicodeFallback{26DD}{\NotoSansSymbolsTwoFont}{^^^^26dd} % Squared Saltire - ⛝
\DeclareUnicodeFallback{26DE}{\NotoSansSymbolsTwoFont}{^^^^26de} % Falling Diagonal In White Circle In Black Square - ⛞
\DeclareUnicodeFallback{26DF}{\NotoSansSymbolsTwoFont}{^^^^26df} % Black Truck - ⛟

\end{minted}
\newpage
\begin{minted}[breaklines, escapeinside=||]{tex}
\DeclareUnicodeFallback{26E0}{\mystixmathfont}{^^^^26e0} % Restricted Left Entry-1 - ⛠
\DeclareUnicodeFallback{26E1}{\mystixmathfont}{^^^^26e1} % Restricted Left Entry-2 - ⛡
\DeclareUnicodeFallback{26E2}{\mystixmathfont}{^^^^26e2} % Astronomical Symbol For Uranus - ⛢
\DeclareUnicodeFallback{26E3}{\mystixmathfont}{^^^^26e3} % Heavy Circle with Stroke and Two Dots Above - ⛣
\DeclareUnicodeFallback{26E4}{\mystixmathfont}{^^^^26e4} % Pentagram - ⛤
\DeclareUnicodeFallback{26E5}{\mystixmathfont}{^^^^26e5} % Right-Handed Interlaced Pentagram - ⛥
\DeclareUnicodeFallback{26E6}{\mystixmathfont}{^^^^26e6} % Left-Handed Interlaced Pentagram - ⛦
\DeclareUnicodeFallback{26E7}{\mystixmathfont}{^^^^26e7} % Inverted Pentagram - ⛧
\DeclareUnicodeFallback{26E8}{\NotoSansSymbolsFont}{^^^^26e8} % Black Cross On Shield - ⛨
\DeclareUnicodeFallback{26E9}{\NotoEmojiFont}{^^^^26e9} % Shinto Shrine - ⛩
\DeclareUnicodeFallback{26EA}{\NotoEmojiFont}{^^^^26ea} % Church - ⛪
\DeclareUnicodeFallback{26EB}{\mystixmathfont}{^^^^26eb} % Castle - ⛫
\DeclareUnicodeFallback{26EC}{\NotoSansSymbolsFont}{^^^^26ec} % Historic Site - ⛬
\DeclareUnicodeFallback{26ED}{\NotoSansSymbolsFont}{^^^^26ed} % Gear Without Hub - ⛭
\DeclareUnicodeFallback{26EE}{\NotoSansSymbolsFont}{^^^^26ee} % Gear with Handles - ⛮
\DeclareUnicodeFallback{26EF}{\NotoSansSymbolsFont}{^^^^26ef} % Map Symbol For Lighthouse - ⛯

\end{minted}
\newpage
\begin{minted}[breaklines, escapeinside=||]{tex}
\DeclareUnicodeFallback{26F0}{\NotoEmojiFont}{^^^^26f0} % Mountain - ⛰
\DeclareUnicodeFallback{26F1}{\NotoEmojiFont}{^^^^26f1} % Umbrella On Ground - ⛱
\DeclareUnicodeFallback{26F2}{\NotoEmojiFont}{^^^^26f2} % Fountain - ⛲
\DeclareUnicodeFallback{26F3}{\NotoEmojiFont}{^^^^26f3} % Flag In Hole - ⛳
\DeclareUnicodeFallback{26F4}{\NotoEmojiFont}{^^^^26f4} % Ferry - ⛴
\DeclareUnicodeFallback{26F5}{\NotoEmojiFont}{^^^^26f5} % Sailboat - ⛵
\DeclareUnicodeFallback{26F6}{\NotoSansSymbolsFont}{^^^^26f6} % Square Four Corners - ⛶
\DeclareUnicodeFallback{26F7}{\NotoEmojiFont}{^^^^26f7} % Skier - ⛷
\DeclareUnicodeFallback{26F8}{\NotoEmojiFont}{^^^^26f8} % Ice Skate - ⛸
\DeclareUnicodeFallback{26F9}{\NotoEmojiFont}{^^^^26f9} % Person with Ball - ⛹
\DeclareUnicodeFallback{26FA}{\NotoEmojiFont}{^^^^26fa} % Tent - ⛺
\DeclareUnicodeFallback{26FB}{\NotoSansSymbolsFont}{^^^^26fb} % Japanese Bank Symbol - ⛻
\DeclareUnicodeFallback{26FC}{\NotoSansSymbolsFont}{^^^^26fc} % Headstone Graveyard Symbol - ⛼
\DeclareUnicodeFallback{26FD}{\NotoEmojiFont}{^^^^26fd} % Fuel Pump - ⛽
\DeclareUnicodeFallback{26FE}{\NotoSansSymbolsFont}{^^^^26fe} % Cup On Black Square - ⛾
\DeclareUnicodeFallback{26FF}{\NotoSansSymbolsFont}{^^^^26ff} % White Flag with Horizontal Middle Black Stripe - ⛿
\end{minted}
\endgroup


\clearpage

\ifinswedish
\printglossary[style=mylong, type=\acronymtype, title={Lista över akronymer
och förkortningar}]
\else
%\printglossary[style=mylong, type=\acronymtype, title={List of Acronyms and abbreviations}]
%\printnoidxglossaries% Replace \printglossary[type=\acronymtype, ...] with:
\printnoidxglossary[type=\acronymtype, style=mylong, title={List of Acronyms and Abbreviations}]
\fi
\clearpage


% --- print the bibliography ---
% Print the bibliography (and make it appear in the table of contents)

% Specify the title of the references page
\renewcommand{\bibname}{References}

%\typeout{Biblatex current language is \currentlang}
\printbibliography[heading=bibintoc, title={References}]


\end{document}

